
\mxunit{stmg-c08}{Révisions et consolidation}
\section{Stratégie de révision}
\begin{keypoint}
Une révision efficace alterne récupération en mémoire, correction explicite des erreurs et réemploi dans un problème. Relire passivement le cours ne suffit pas : il faut tenter une question sans modèle, comparer sa réponse, puis refaire la question après un délai.
\end{keypoint}

Pour chaque chapitre, utiliser le code suivant : \textbf{V} (maîtrisé), \textbf{O} (fragile), \textbf{R} (à reprendre). Après le diagnostic, travailler en priorité une compétence orange ou rouge, puis vérifier immédiatement avec une question analogue.

\begin{center}
\begin{tabularx}{.9\textwidth}{lcccX}
\toprule
Compétence & V & O & R & Action de remédiation \\
\midrule
Choisir un modèle & $\square$ & $\square$ & $\square$ & Refaire un exercice de tri arithmétique/géométrique/exponentiel.\\
Lire et calculer & $\square$ & $\square$ & $\square$ & Reprendre les automatismes ciblés.\\
Rédiger et interpréter & $\square$ & $\square$ & $\square$ & Ajouter données, formule, calcul, unité et phrase de réponse.\\
\bottomrule
\end{tabularx}
\end{center}

\mxsession{52}{\mxdate*{2027-05-11}}{B : 2 h -- C : 1 h}{Diagnostic des acquis}{Révisions et consolidation}
\begin{objectives}\begin{itemize}\item Identifier les acquis.\item Classer les erreurs.\item Planifier la remédiation.\end{itemize}\end{objectives}
\begin{fluency}[6 à 8 minutes, sans calculatrice]\begin{enumerate}\item Écrire le coefficient multiplicateur d'une hausse de $\pc{25}$. \mxblank{$1+\dfrac{25}{100}=\num{1,25}$}\item Donner la formule du terme général arithmétique depuis $u_0$. \mxblank{$u_n=u_0+nr$}\item Résoudre $10^x=1000$. \mxblank{$x=3$}\item Dériver $x+\dfrac4x$. \mxblank{$1-\dfrac4{x^2}$}\item Énoncer la formule des probabilités totales avec deux branches. \mxblank{$P(B)=P(A)P_A(B)+P(\overline A)P_{\overline A}(B)$}\end{enumerate}\end{fluency}
\begin{prerequisites}Relire la dernière trace écrite et savoir effectuer les calculs élémentaires mobilisés aujourd'hui. Avant l'activité, faire reformuler les mots-clés : identifier les acquis et classer les erreurs.\end{prerequisites}

\begin{activity}
Répondre sans cours à 14 questions, deux par chapitre. Reporter chaque erreur dans une grille : connaissance, choix de méthode, calcul, lecture, rédaction. Choisir ensuite deux priorités personnelles.

\end{activity}
\begin{solution}La correction attend un profil d'erreurs et un plan de reprise, pas seulement une note globale. Faire verbaliser la donnée qui déclenche le choix de méthode.\end{solution}

\begin{summary}Le diagnostic sert à choisir la remédiation : reconnaître l'erreur, retrouver la règle utile, refaire une question analogue, puis vérifier sans aide.\end{summary}

\begin{example}
Reprendre les données de l'activité et rédiger la solution dans l'ordre : définition des objets, propriété utilisée, calcul exact ou approché, puis phrase de réponse. La réponse modèle figure dans l'édition enseignant.
\end{example}
\begin{solution}La correction attend un profil d'erreurs et un plan de reprise, pas seulement une note globale.\end{solution}

\begin{exercise}[essentiel]
\begin{enumerate}\item Choisir entre suite arithmétique, suite géométrique et fonction exponentielle dans trois situations.\item Résoudre une équation exponentielle avec un logarithme.\item Calculer une probabilité à partir d'un arbre puis interpréter.\end{enumerate}
\end{exercise}
\begin{solution}Justifier chaque choix par variation absolue ou relative ; vérifier la positivité avant le logarithme ; multiplier le long d'un chemin et conclure par une phrase.\end{solution}

\mxexuse{52}

\begin{homework}Reprendre l'exercice d'entraînement 3 avec les données suivantes : remplacer la donnée initiale par $380$ et le taux ou la raison par $6$. Rédiger en quatre lignes : données, formule, calcul, conclusion.\end{homework}
\begin{solution}La correction dépend de l'exercice repris. Vérifier surtout la cohérence du modèle, le nombre d'itérations, l'arrondi et la phrase de conclusion.\end{solution}

\begin{exitticket}[3 minutes]Sur une demi-feuille : écrire la règle principale de la séance, traiter un exemple très court et indiquer son degré de confiance de 1 à 4. Question(s) essentielle(s) du jour : 1 et 2.\end{exitticket}

\ifmxteacher
\begin{method}
\textbf{Conduite en 1 h :} 7 min d'automatismes ; 10 min de mise en commun et trace ; 28 min d'exercices essentiels ; 5 min de correction ; 3 min de billet de sortie.\\
\textbf{Conduite en 2 h :} ajouter 20 min de recherche de l'activité, 20 min d'institutionnalisation, un second exercice autonome et 10 min de correction comparée.\\
\textbf{Différenciation :} fournir aux élèves fragiles une fiche à quatre cases (données, modèle, calcul, conclusion). Demander aux élèves rapides de produire un contre-exemple, une justification ou une seconde méthode.
\end{method}
\fi
\mxsession{53}{\mxdate*{2027-05-13}}{B : 1 h -- C : 2 h}{Réactivation : fonctions exponentielles}{Révisions et consolidation}
\begin{objectives}\begin{itemize}\item Réactiver les modèles exponentiels.\item Résoudre et interpréter.\item Contrôler les unités.\end{itemize}\end{objectives}
\begin{fluency}[6 à 8 minutes, sans calculatrice]\begin{enumerate}\item Réduire $3x+7-x+2$. \mxblank{$2x+9$}\item Résoudre $10^x=1000$. \mxblank{$x=3$}\item Dériver $x+\dfrac4x$. \mxblank{$1-\dfrac4{x^2}$}\item Énoncer la formule des probabilités totales avec deux branches. \mxblank{$P(B)=P(A)P_A(B)+P(\overline A)P_{\overline A}(B)$}\item Donner la formule du terme général arithmétique depuis $u_0$. \mxblank{$u_n=u_0+nr$}\end{enumerate}\end{fluency}
\begin{prerequisites}Relire la dernière trace écrite et savoir effectuer les calculs élémentaires mobilisés aujourd'hui. Avant l'activité, faire reformuler les mots-clés : réactiver les modèles exponentiels et résoudre et interpréter.\end{prerequisites}

\begin{activity}
Une valeur de $\num{18000}$ augmente de $\pc{3,5}$ par an. Écrire le modèle, calculer la valeur après huit ans et déterminer quand elle dépassera $\num{25000}$.
\begin{center}\begin{tikzpicture}\begin{axis}[width=.68\textwidth,height=5cm,axis lines=middle,xmin=-3,xmax=4,ymin=0,ymax=5,grid=major]\addplot[mxteacher,very thick,samples=100,domain=-3:4]{1.4^x};\addplot[mxwarn,very thick,samples=100,domain=-3:4]{.7^x};\end{axis}\end{tikzpicture}\end{center}
\end{activity}
\begin{solution}$V(t)=\num{18000}\times\num{1,035}^t$, $V(8)\approx\num{23703}$. Le seuil est dépassé à partir de la dixième année. Faire verbaliser la donnée qui déclenche le choix de méthode.\end{solution}

\begin{summary}Une fonction exponentielle modélise une évolution relative constante. Identifier d'abord la valeur initiale et le coefficient multiplicateur.\end{summary}

\begin{example}
Reprendre les données de l'activité et rédiger la solution dans l'ordre : définition des objets, propriété utilisée, calcul exact ou approché, puis phrase de réponse. La réponse modèle figure dans l'édition enseignant.
\end{example}
\begin{solution}$V(t)=\num{18000}\times\num{1,035}^t$, $V(8)\approx\num{23703}$. Le seuil est dépassé à partir de la dixième année.\end{solution}

\begin{exercise}[essentiel]
\begin{enumerate}\item Choisir entre suite arithmétique, suite géométrique et fonction exponentielle dans trois situations.\item Résoudre une équation exponentielle avec un logarithme.\item Calculer une probabilité à partir d'un arbre puis interpréter.\end{enumerate}
\end{exercise}
\begin{solution}Justifier chaque choix par variation absolue ou relative ; vérifier la positivité avant le logarithme ; multiplier le long d'un chemin et conclure par une phrase.\end{solution}

\mxexuse{53}

\begin{homework}Reprendre l'exercice d'entraînement 2 avec les données suivantes : remplacer la donnée initiale par $385$ et le taux ou la raison par $7$. Rédiger en quatre lignes : données, formule, calcul, conclusion.\end{homework}
\begin{solution}La correction dépend de l'exercice repris. Vérifier surtout la cohérence du modèle, le nombre d'itérations, l'arrondi et la phrase de conclusion.\end{solution}

\begin{exitticket}[3 minutes]Sur une demi-feuille : écrire la règle principale de la séance, traiter un exemple très court et indiquer son degré de confiance de 1 à 4. Question(s) essentielle(s) du jour : 1, 2 et 3.\end{exitticket}

\ifmxteacher
\begin{method}
\textbf{Conduite en 1 h :} 7 min d'automatismes ; 10 min de mise en commun et trace ; 28 min d'exercices essentiels ; 5 min de correction ; 3 min de billet de sortie.\\
\textbf{Conduite en 2 h :} ajouter 20 min de recherche de l'activité, 20 min d'institutionnalisation, un second exercice autonome et 10 min de correction comparée.\\
\textbf{Différenciation :} fournir aux élèves fragiles une fiche à quatre cases (données, modèle, calcul, conclusion). Demander aux élèves rapides de produire un contre-exemple, une justification ou une seconde méthode.
\end{method}
\fi
\mxsession{54}{\mxdate*{2027-05-18}}{B : 2 h -- C : 1 h}{Réactivation : logarithme décimal}{Révisions et consolidation}
\begin{objectives}\begin{itemize}\item Réactiver les logarithmes.\item Résoudre une équation.\item Vérifier les conditions.\end{itemize}\end{objectives}
\begin{fluency}[6 à 8 minutes, sans calculatrice]\begin{enumerate}\item Calculer $\pc{5}$ de $160$. \mxblank{$8$}\item Dériver $x+\dfrac4x$. \mxblank{$1-\dfrac4{x^2}$}\item Énoncer la formule des probabilités totales avec deux branches. \mxblank{$P(B)=P(A)P_A(B)+P(\overline A)P_{\overline A}(B)$}\item Donner la formule du terme général arithmétique depuis $u_0$. \mxblank{$u_n=u_0+nr$}\item Résoudre $10^x=1000$. \mxblank{$x=3$}\end{enumerate}\end{fluency}
\begin{prerequisites}Relire la dernière trace écrite et savoir effectuer les calculs élémentaires mobilisés aujourd'hui. Avant l'activité, faire reformuler les mots-clés : réactiver les logarithmes et résoudre une équation.\end{prerequisites}

\begin{activity}
Résoudre $10^x=340$, $\num{1,06}^n=2$ et $\num{1,06}^n\geq2$. Pour chaque réponse, distinguer valeur réelle, valeur approchée et rang entier minimal.

\end{activity}
\begin{solution}$x\approx\num{2,5315}$ ; $n\approx\num{11,90}$ ; pour le seuil entier, $n=12$. Faire verbaliser la donnée qui déclenche le choix de méthode.\end{solution}

\begin{summary}Le logarithme transforme une équation exponentielle en équation sur l'exposant. Conditions : base positive différente de $1$, second membre positif.\end{summary}

\begin{example}
Reprendre les données de l'activité et rédiger la solution dans l'ordre : définition des objets, propriété utilisée, calcul exact ou approché, puis phrase de réponse. La réponse modèle figure dans l'édition enseignant.
\end{example}
\begin{solution}$x\approx\num{2,5315}$ ; $n\approx\num{11,90}$ ; pour le seuil entier, $n=12$.\end{solution}

\begin{exercise}[essentiel]
\begin{enumerate}\item Choisir entre suite arithmétique, suite géométrique et fonction exponentielle dans trois situations.\item Résoudre une équation exponentielle avec un logarithme.\item Calculer une probabilité à partir d'un arbre puis interpréter.\end{enumerate}
\end{exercise}
\begin{solution}Justifier chaque choix par variation absolue ou relative ; vérifier la positivité avant le logarithme ; multiplier le long d'un chemin et conclure par une phrase.\end{solution}

\mxexuse{54}

\begin{homework}Reprendre l'exercice d'entraînement 3 avec les données suivantes : remplacer la donnée initiale par $390$ et le taux ou la raison par $2$. Rédiger en quatre lignes : données, formule, calcul, conclusion.\end{homework}
\begin{solution}La correction dépend de l'exercice repris. Vérifier surtout la cohérence du modèle, le nombre d'itérations, l'arrondi et la phrase de conclusion.\end{solution}

\begin{exitticket}[3 minutes]Sur une demi-feuille : écrire la règle principale de la séance, traiter un exemple très court et indiquer son degré de confiance de 1 à 4. Question(s) essentielle(s) du jour : 1 et 2.\end{exitticket}

\ifmxteacher
\begin{method}
\textbf{Conduite en 1 h :} 7 min d'automatismes ; 10 min de mise en commun et trace ; 28 min d'exercices essentiels ; 5 min de correction ; 3 min de billet de sortie.\\
\textbf{Conduite en 2 h :} ajouter 20 min de recherche de l'activité, 20 min d'institutionnalisation, un second exercice autonome et 10 min de correction comparée.\\
\textbf{Différenciation :} fournir aux élèves fragiles une fiche à quatre cases (données, modèle, calcul, conclusion). Demander aux élèves rapides de produire un contre-exemple, une justification ou une seconde méthode.
\end{method}
\fi
\mxsession{55}{\mxdate*{2027-05-20}}{B : 1 h -- C : 2 h}{Réactivation : fonction inverse}{Révisions et consolidation}
\begin{objectives}\begin{itemize}\item Réactiver la fonction inverse.\item Mener une étude.\item Tracer et interpréter.\end{itemize}\end{objectives}
\begin{fluency}[6 à 8 minutes, sans calculatrice]\begin{enumerate}\item Écrire le coefficient multiplicateur d'une hausse de $\pc{10}$. \mxblank{$1+\dfrac{10}{100}=\num{1,10}$}\item Énoncer la formule des probabilités totales avec deux branches. \mxblank{$P(B)=P(A)P_A(B)+P(\overline A)P_{\overline A}(B)$}\item Donner la formule du terme général arithmétique depuis $u_0$. \mxblank{$u_n=u_0+nr$}\item Résoudre $10^x=1000$. \mxblank{$x=3$}\item Dériver $x+\dfrac4x$. \mxblank{$1-\dfrac4{x^2}$}\end{enumerate}\end{fluency}
\begin{prerequisites}Relire la dernière trace écrite et savoir effectuer les calculs élémentaires mobilisés aujourd'hui. Avant l'activité, faire reformuler les mots-clés : réactiver la fonction inverse et mener une étude.\end{prerequisites}

\begin{activity}
Étudier $f(x)=x+\dfrac{16}{x}$ sur $]0,+\infty[$ et interpréter son minimum dans un contexte de coût moyen. Produire le tableau de variations et une esquisse.
\begin{center}\begin{tikzpicture}\begin{axis}[width=.68\textwidth,height=5cm,axis lines=middle,xmin=-4,xmax=4,ymin=-4,ymax=4,grid=major]\addplot[mxteacher,very thick,samples=80,domain=-4:-.15]{1/x};\addplot[mxteacher,very thick,samples=80,domain=.15:4]{1/x};\end{axis}\end{tikzpicture}\end{center}
\end{activity}
\begin{solution}$f'(x)=\dfrac{(x-4)(x+4)}{x^2}$ ; minimum $f(4)=8$. Faire verbaliser la donnée qui déclenche le choix de méthode.\end{solution}

\begin{summary}Pour une fonction contenant $\dfrac1x$, commencer par exclure $x=0$, puis dériver, factoriser et étudier le signe avant les variations.\end{summary}

\begin{example}
Reprendre les données de l'activité et rédiger la solution dans l'ordre : définition des objets, propriété utilisée, calcul exact ou approché, puis phrase de réponse. La réponse modèle figure dans l'édition enseignant.
\end{example}
\begin{solution}$f'(x)=\dfrac{(x-4)(x+4)}{x^2}$ ; minimum $f(4)=8$.\end{solution}

\begin{exercise}[essentiel]
\begin{enumerate}\item Choisir entre suite arithmétique, suite géométrique et fonction exponentielle dans trois situations.\item Résoudre une équation exponentielle avec un logarithme.\item Calculer une probabilité à partir d'un arbre puis interpréter.\end{enumerate}
\end{exercise}
\begin{solution}Justifier chaque choix par variation absolue ou relative ; vérifier la positivité avant le logarithme ; multiplier le long d'un chemin et conclure par une phrase.\end{solution}

\mxexuse{55}

\begin{homework}Reprendre l'exercice d'entraînement 2 avec les données suivantes : remplacer la donnée initiale par $395$ et le taux ou la raison par $3$. Rédiger en quatre lignes : données, formule, calcul, conclusion.\end{homework}
\begin{solution}La correction dépend de l'exercice repris. Vérifier surtout la cohérence du modèle, le nombre d'itérations, l'arrondi et la phrase de conclusion.\end{solution}

\begin{exitticket}[3 minutes]Sur une demi-feuille : écrire la règle principale de la séance, traiter un exemple très court et indiquer son degré de confiance de 1 à 4. Question(s) essentielle(s) du jour : 1, 2 et 3.\end{exitticket}

\ifmxteacher
\begin{method}
\textbf{Conduite en 1 h :} 7 min d'automatismes ; 10 min de mise en commun et trace ; 28 min d'exercices essentiels ; 5 min de correction ; 3 min de billet de sortie.\\
\textbf{Conduite en 2 h :} ajouter 20 min de recherche de l'activité, 20 min d'institutionnalisation, un second exercice autonome et 10 min de correction comparée.\\
\textbf{Différenciation :} fournir aux élèves fragiles une fiche à quatre cases (données, modèle, calcul, conclusion). Demander aux élèves rapides de produire un contre-exemple, une justification ou une seconde méthode.
\end{method}
\fi
